% Modernised reading — fonds Alexandre Grothendieck, shelfmark 29, batch 9
% (pages 161-180).
%
% This is an INTERPRETATION, not a transcription. It re-reads the pages in
% today's notation and vocabulary, states the hypotheses the manuscript leaves
% implicit, and drops the critical apparatus so that the mathematics reads
% straight through. Everything the brackets and underlines carried moves into a
% footnote. In any dispute of substance the transcription governs: whoever
% cites Grothendieck must cite batch-09.fr.tex.
%
% Scope: the folder was read WHOLE before any of this was written — all eleven
% batches, pages 1-216, in one pass — and one file was then written per batch.
% The one explanation this file could not have given batch by batch is why the
% open question of page 179 is open exactly in the case page 170 calls the
% interesting one: for G finite the constant group scheme is affine over S' and
% Grothendieck's own descent theorem settles effectivity; for G infinite it is
% not even quasi-compact, and the theorem does not reach it.
%
% Pages 164-169 are the hardest run in the folder: the formulas carry, the
% sentences joining them very often do not, and the lower third of page 166 is
% cancelled and unreadable. This reading states the statements and marks the
% joins it cannot supply.
%
% Departures from the page, each footnoted where it occurs:
%   - page 161's nodal counter-example: pi_1(Y) is the FULL profinite
%     completion of Z, not Z-hat deprived of its p-component. The cyclic
%     necklace coverings exist in every degree, p included, because they come
%     from the dual graph and not from the geometry. The page's restriction is
%     an uncertain reading, and it is unnecessary either way;
%   - the same example needs X = P^1 for pi_1(X) = 0;
%   - « conoyau fini » for a map of profinite groups is written here as « image
%     d'indice fini », which is what the statement means;
%   - page 178 writes « contravariant » and strikes it: the functor is
%     covariant, and the reading says so;
%   - page 175 writes the exponent as a roman I where the rest of the page
%     writes pi^2.
%
% Pass: Opus 5 (claude-opus-5), 2026-08-29 — first pass, unchecked against the
% pages by a human.
\documentclass[11pt,a4paper]{article}
% Le préambule commun à toutes les transcriptions du fonds Grothendieck.
%
% Il tient en une idée : **l'incertitude se compose, elle ne se résout pas.**
% Une transcription qui lisse un mot douteux a détruit la seule chose pour
% laquelle on la faisait — un lecteur doit pouvoir savoir, à chaque endroit,
% ce qui était sur le feuillet et ce qui a été ajouté. Les six macros qui
% suivent sont donc l'appareil critique, et non de la décoration.
%
% Les mêmes commandes sont reconnues par `scripts/render.mjs`, qui produit la
% vue de lecture du site. Une macro ajoutée ici doit l'être aussi là-bas,
% faute de quoi le rendu échouera bruyamment — ce qui est voulu : un rendu
% silencieusement faux est pire qu'un rendu absent.

\ProvidesPackage{grothendieck}[2026/08/08 Transcriptions du fonds Grothendieck]

\usepackage{amsmath,amssymb,amsthm}
\usepackage{mathrsfs}
\usepackage{stmaryrd}
% Les diagrammes commutatifs sont partout dans ce fonds : c'est la langue
% même de la géométrie algébrique de Grothendieck, et une transcription qui
% les rendrait en prose ne transcrirait rien.
\usepackage{tikz-cd}
% Français par défaut — c'est la langue des feuillets — mais l'anglais est
% chargé aussi : la traduction et le résumé sont des documents anglais, et la
% typographie française (espaces avant les deux-points) y serait une faute.
% Ces documents basculent par \selectlanguage{english} en tête de corps.
\usepackage[english,french]{babel}
\usepackage{fontspec}
\usepackage{geometry}
\geometry{a4paper,margin=25mm,marginparwidth=28mm}
\usepackage{marginnote}
\usepackage{xcolor}
% ulem, not soul. `soul`'s \st and \ul are built on a character-by-character
% scan that XeLaTeX's font machinery breaks: the document fails with an
% undefined control sequence rather than degrading. ulem does the same two jobs
% and is Unicode-engine safe. `normalem` stops it hijacking \emph, which the
% transcriptions use for Grothendieck's own emphasis.
\usepackage[normalem]{ulem}
\usepackage{hyperref}
\hypersetup{colorlinks=true,linkcolor=black,urlcolor=black,pdfborder={0 0 0}}

% --- Métadonnées du lot -----------------------------------------------------
% Reprises telles quelles par le script de rendu, qui les affiche en tête de
% la vue de lecture. Elles fixent ce que le fichier prétend couvrir : sans
% elles, une transcription flotte sans point d'attache dans un fonds de seize
% mille feuillets.

\newcommand{\folder}[1]{\gdef\grothFolder{#1}}
\newcommand{\batch}[1]{\gdef\grothBatch{#1}}
\newcommand{\leaves}[2]{\gdef\grothFirst{#1}\gdef\grothLast{#2}}
\newcommand{\foldertitle}[1]{\gdef\grothFoldertitle{#1}}
\gdef\grothFolder{?}\gdef\grothBatch{?}\gdef\grothFirst{?}\gdef\grothLast{?}\gdef\grothFoldertitle{}

% --- Le résumé --------------------------------------------------------------
%
% Il ouvre la lecture modernisée, sous le titre « Résumé », et il est composé
% en petit. Ce n'est pas un ornement : c'est le seul passage du document qui
% ne soit pas des mathématiques, et un lecteur doit pouvoir le reconnaître
% avant de l'avoir lu — puis le sauter s'il connaît déjà le sujet, ou s'y
% arrêter s'il ne le connaît pas. Le filet qui le suit marque la reprise du
% corps.

\newenvironment{resume}
  {\small\setlength{\parskip}{0.45em}}
  {\par\medskip\hrule\medskip}

% --- Mots-clés ---------------------------------------------------------------
%
% En anglais, en clôture du résumé de la lecture modernisée : le vocabulaire
% moderne sous lequel chercher ce que les feuillets construisent. Le manifeste
% du site (`npm run manifest`) les extrait pour étiqueter la cote entière —
% cette ligne est la source unique des tags, il n'y a pas de fichier à côté.
\newcommand{\keywords}[1]{\par\smallskip\noindent
  {\footnotesize\textsc{Keywords} --- \textit{#1}}\par}

% --- L'appareil critique ----------------------------------------------------

% Le numéro de feuillet, en marge. C'est l'ancre qui permet de régler un
% désaccord sur une formule en regardant *un* feuillet plutôt qu'une chemise
% de six cents. Le site s'en sert aussi pour tourner les pages du fac-similé
% quand on fait défiler la transcription.
\newcommand{\leaf}[1]{%
  \marginnote{\footnotesize\color{gray!70}\textsf{#1}}%
  \label{leaf:#1}\ignorespaces}

% Illisible : on ne devine pas. Un mot inventé qui se lit comme les autres est
% le pire résultat possible.
\newcommand{\ill}{\textcolor{red!60!black}{[\,\ldots\,]}}

% Lecture douteuse : le mot est donné, et signalé comme incertain.
\newcommand{\uncertain}[1]{\uline{#1}}

% Ajout de l'éditeur — une lettre manquante, un mot sous-entendu.
\newcommand{\add}[1]{\textcolor{blue!50!black}{[#1]}}

% Biffé par Grothendieck lui-même. Ce qu'il a rayé fait partie du document :
% une rature dit souvent où la pensée a changé de direction.
\newcommand{\struck}[1]{\sout{#1}}

% Note du transcripteur, sur l'état matériel du feuillet.
\newcommand{\note}[1]{\par\smallskip{\small\color{gray!60!black}\textit{[#1]}}\par\smallskip}

% Note marginale de Grothendieck, distincte du corps du texte.
\newcommand{\marginal}[1]{\par\smallskip\noindent\hspace{1em}%
  {\small\color{gray!60!black}#1}\par\smallskip}

% --- Titre ------------------------------------------------------------------

\AtBeginDocument{%
  \begin{center}
    {\large\bfseries Fonds Alexandre Grothendieck --- cote n\textsuperscript{o}~\grothFolder}\\[2pt]
    {\itshape\grothFoldertitle}\\[4pt]
    {\small feuillets \grothFirst--\grothLast\ (lot \grothBatch)}\\[2pt]
    {\footnotesize\color{gray!60!black}Transcription établie d'après le fac-similé numérisé
     par l'Université de Montpellier.\\
     Les lectures incertaines sont soulignées, les illisibles notées [\,\ldots\,],
     les ajouts entre crochets.}
  \end{center}
  \vspace{1em}}

\endinput


\folder{29}
\batch{9}
\pages{161}{180}
\dating{1959-1969}
\watermark{Édition de démonstration\\interprétation personnelle de l'œuvre}
\foldertitle{Groupe fondamental [Autour de SGA 4 et SGA 7] — lecture modernisée}

\begin{document}

\section*{Résumé}

\begin{resume}
Jusqu'ici, dans ce carton, un revêtement était fini. Le groupe fondamental
classait des ensembles à un nombre fini d'éléments, et c'est pourquoi il était
\emph{profini} : une limite de groupes finis. Ces pages demandent ce qui se
passe quand on lève cette restriction.

La question n'est pas gratuite. Beaucoup d'objets naturels ne sont pas finis
mais sont localement constants : un système local de $\mathbb{Z}$-modules, un
faisceau de groupes abéliens de type fini, un schéma en groupes commutatifs
sans torsion. Ils devraient être classés par un groupe fondamental, mais celui
qu'on a ne les voit pas. Grothendieck écrit la phrase qui justifie tout le
lot : « Notre intérêt ici est le cas $G$ \emph{infini} ».

Le procédé est celui de la lettre à Serre de 1959 --- construire le foncteur
plutôt que l'objet ---, mais il faut cette fois travailler. Le foncteur qui, à
un groupe $\pi$, associe l'ensemble des fibrés principaux de groupe $\pi$ munis
d'un point marqué, est \emph{exact à gauche}, et Grothendieck le démontre en
trois temps ; puis, sous une hypothèse noethérienne, il montre que les groupes
auxquels on peut restreindre le groupe structural forment un système
\emph{décroissant} filtrant, dont l'intersection porte encore un fibré. D'où la
pro-représentabilité, et un groupe fondamental élargi, qui classifie
précisément les objets localement constants à fibres infinies mais de type
fini.

Deux autres choses valent d'être vues.

La première est un contre-exemple, à la page 161, et c'est la seule page du
carton où Grothendieck démontre qu'un théorème est faux. Le théorème dit que
pour un morphisme propre dominant vers une base \emph{normale}, l'image des
groupes fondamentaux est d'indice fini. Ôtez la normalité : prenez pour base
une courbe rationnelle à point double, et pour source sa normalisée. La source
est une droite projective, dont le groupe fondamental est trivial ; la base,
elle, a des revêtements en collier --- $n$ droites recollées en cycle --- de
tout degré. L'image est donc triviale dans un groupe infini. La normalité n'est
pas une commodité, c'est l'hypothèse.

La seconde est un aveu, à la page 179, et c'est le second du carton après celui
de la page 115. Grothendieck démontre que les données de descente
\emph{effectives} forment aussi un foncteur pro-représentable, puis écrit :
« J'ignore \dots\ si dans ce cas $\pi^{1} = \pi^{1}_{\mathrm{eff}}$, i.e.
\emph{toute} donnée de descente sur $G_{S'}$ est \emph{effective} ». La raison
pour laquelle il l'ignore est instructive et elle se lit en rapprochant deux
pages : quand $G$ est fini, le faisceau constant $G_{S'}$ est fini donc affine,
et le théorème de descente de Grothendieck lui-même règle la question ; quand
$G$ est infini, $G_{S'}$ est une somme infinie de copies de $S'$, donc pas même
quasi-compact, et le théorème ne l'atteint plus. La question est ouverte
exactement dans le cas qui l'intéresse.

Les noms modernes à chercher sont \emph{groupe fondamental élargi},
\emph{groupe fondamental pro-étale}, \emph{pro-représentabilité},
\emph{descente effective}, \emph{homotopie étale}.

\keywords{pro-etale fundamental group, enlarged fundamental group, strict pro-representability, effective descent datum, etale homotopy}
\end{resume}

\pagerange{161}{162}
\section*{La normalité n'est pas une commodité (pages 161 et 162)}

\subsection*{L'énoncé}

Soit $f : X \to Y$ un morphisme propre dominant, $Y$ \emph{normal} connexe et
localement noethérien. Alors l'image de
\[
  \pi_{1}(X) \longrightarrow \pi_{1}(Y)
\]
est d'indice fini.\footnote{La page écrit « a un conoyau fini ». Pour un
morphisme de groupes profinis l'image n'est pas nécessairement distinguée et il
n'y a pas de conoyau ; ce que l'énoncé veut dire, et ce que la démonstration
donne, est que l'image est d'\emph{indice fini}. C'est la formulation retenue
ici.} La démonstration se ramène au cas $Y = \operatorname{Spec}(K)$, en
remplaçant $K$ par une extension séparable finie convenable et $X$ par une
composante connexe de $X'$.

\subsection*{Le contre-exemple}

C'est la moitié lisible de la page, et c'en est le propos :

\begin{quote}
Si $Y$ n'est pas \emph{normal}, le théorème est faux.
\end{quote}

Prenons pour $Y$ une courbe rationnelle sur $k$ algébriquement clos ayant un
point double ordinaire --- une cubique nodale, par exemple --- et pour $X$ sa
normalisée. Alors $X \simeq \mathbb{P}^{1}_{k}$, donc
\[
  \pi_{1}(X) \;=\; 0 ,
\]
tandis que
\[
  \pi_{1}(Y) \;\simeq\; \hat{\mathbb{Z}} .
\]
L'image est triviale dans un groupe infini : elle n'est pas d'indice fini.

Le groupe $\pi_{1}(Y)$ se voit à l'œil. Prenons $n$ copies
$C_{1}, \dots, C_{n}$ de $\mathbb{P}^{1}$ et recollons $\infty_{i}$ à
$0_{i+1}$ cycliquement ; le collier ainsi obtenu s'envoie sur $Y$, chaque
$C_{i}$ par la normalisation, et le morphisme est fini étale de degré $n$ ---
au point double, les deux branches s'envoient isomorphiquement sur les deux
branches. C'est un revêtement galoisien de groupe $\mathbb{Z}/n$, et il existe
\emph{pour tout $n$}. Le groupe fondamental est donc $\hat{\mathbb{Z}}$ tout
entier.\footnote{La page porte « $\pi_{1}(Y) = \hat{\mathbb{Z}}$
privé de sa $p$-composante », ces derniers mots étant une lecture que la
transcription signale comme douteuse. Elle n'est en tout cas pas nécessaire :
les revêtements en collier existent en degré $p$ comme en tout autre degré,
parce qu'ils viennent du graphe dual de la courbe et non de sa géométrie ---
c'est le $H^{1}$ d'un cercle, qui ne connaît pas la caractéristique. La
restriction au premier à $p$ est celle qui s'impose pour les revêtements
\emph{d'une courbe lisse}, où le $p$-groupe fondamental est sauvage ; elle ne
s'impose pas ici. Le contre-exemple est d'ailleurs plus fort sans elle.}
Enfin, l'hypothèse $X \simeq \mathbb{P}^{1}$ --- c'est-à-dire $Y$
\emph{rationnelle} --- est ce qui donne $\pi_{1}(X) = 0$ ; la page ne l'écrit
pas et elle est nécessaire.

Grothendieck ajoute une observation qui explique où le raisonnement casserait :
dans ce cas, on ne peut pas former $\underline{\mathrm{Hom}}_{Y}(X, Y \sqcup
Y)$. C'est le schéma des morphismes du lemme 1 de la page 47, et son absence
est ce qui interdit de descendre l'isomorphisme générique. La normalité de $Y$
était ce qui le faisait exister.

La page 162 ne porte qu'un carré encadré, $f_{*}(X_{1})$ au-dessus de
$S \leftarrow T$, et rien d'autre.

\pagerange{164}{169}
\section*{Le progroupe fondamental pour les revêtements infinis (pages 164 à 169)}

Le feuillet qui ouvre la suite porte de sa main, avec une rature :
« \emph{Le Groupe} \; Progroupe fondamental pour les revêtements infinis
principaux ». La rature est où le nom a changé : ce n'est pas un groupe, c'est
un pro-groupe.

\subsection*{Le foncteur}

Soit $S$ un préschéma connexe et $\xi$ un point géométrique. Pour tout groupe
$\pi$ --- \emph{pas nécessairement fini} --- on note $\pi^{1}(S, \xi\,;\, \pi)$ l'ensemble des classes d'isomorphie de
fibrés principaux homogènes sur $S$, de groupe structural $\pi$, pointés
au-dessus de $\xi$.

C'est le même foncteur que le $Z(S,a;G)$ de la lettre à Serre et que le
$\mathfrak{Z}(S,a;G)$ des feuilles anciennes, avec la restriction de finitude
levée ; et c'est cette levée qui fait toute la difficulté du lot.

\subsection*{$\alpha$) Le foncteur est exact à gauche}

\textbf{a) Il commute aux produits finis.} Trivial.

\textbf{b) Il transforme une injection en une injection.} Soit
$\pi \hookrightarrow \pi'$. La raison est le lemme d'unicité déjà rencontré :
entre deux $\pi'$-fibrés pointés il existe \emph{au plus un} isomorphisme
respectant les points marqués. Si $P'_{1}$ et $P'_{2}$ proviennent de
$\pi$-fibrés $P_{i} \subset P'_{i}$ avec $\xi_{i} = \xi'_{i}$, l'unique
isomorphisme $u' : P'_{1} \to P'_{2}$ se restreint donc en
$u : P_{1} \to P_{2}$.

Et la page 165 dit exactement d'où vient la réduction du groupe structural :
\emph{$P_{i}$ se déduit de $P'_{i}$ par une section de $P'_{i}/\pi$ pointée par
$\xi'_{i}$ ; et une telle section, si elle existe, est unique}. Réduire le
groupe structural, c'est trancher une section ; le point marqué la rend unique.

\textbf{c) Il préserve l'exactitude d'un couple.} Si
$\pi \to \pi' \rightrightarrows \pi''$ est exacte, il en va de même des images.
On considère $Q$, le sous-schéma des coïncidences des deux morphismes
$u_{1}, u_{2} : P' \to P''$ ; il est ouvert et fermé, contient $\xi'$, est
stable par $\pi$, et lui correspond donc un sous-préschéma ouvert et fermé
$S'$ de $P'/\pi$, fidèlement plat et quasi-compact sur
$S$.\footnote{Les pages 164 et 165 sont d'une écriture très rapide et la
transcription y laisse la plupart des phrases de liaison illisibles. Les trois
énoncés a), b), c) et les deux raisons portantes --- l'unicité de l'isomorphisme
pointé, l'unicité de la section --- sont ce que la page donne ; le détail des
vérifications ne l'est pas et n'est pas reconstruit ici.}

\subsection*{d) La filtration décroissante}

C'est le point nouveau, et il demande une hypothèse : \emph{$S$ localement
noethérien, ou bien $\pi$ artinien} (condition de chaîne descendante).

Soit $c \in \pi^{1}(S,\xi;\pi)$ correspondant à $(P, \xi)$. Considérons deux
sous-groupes $\pi'$, $\pi''$ auxquels on peut restreindre le groupe structural.
Alors
\[
  P/(\pi' \cap \pi'') \longrightarrow P/\pi' \times P/\pi''
\]
est une \emph{immersion ouverte et fermée}, et le point $(y', y'')$ donné par
les deux réductions est dans son image. On peut donc restreindre le groupe
structural à $\pi' \cap \pi''$.

Il en résulte que \emph{les groupes auxquels on peut restreindre forment un
système filtrant décroissant} $(\pi_{i})$, avec des $P_{i} \subset P$
correspondants. C'est l'analogue, pour des groupes infinis, de la condition de
minimalité de la lettre de 1959 : là, la chaîne descendante était stationnaire
par hypothèse ; ici elle ne l'est pas, et il faut montrer que l'intersection
porte encore un fibré.

\subsection*{L'intersection porte un fibré}

C'est ce que démontrent les pages 167 et 169.

La question est locale sur $S$, et par étalement sur un $S^{(0)}$ de type fini
sur $\mathbb{Z}$ on peut supposer $S$ \emph{noethérien}. Soit alors $P^{0}$ la
composante connexe de $\xi$ dans $P$ : elle est ouverte et fermée, $P$ étant
localement noethérien donc localement connexe.

\textbf{Assertion.} $P^{0} \to S$ est \emph{surjectif}.

L'image est constructible et stable par générisation ; on se ramène au cas d'un
anneau de valuation complet, où l'on conclut.

Alors, puisque $P^{0} \subset \bigcap_{i} P^{(i)}$ et que chacun des $P^{(i)}$
est une réunion de composantes connexes de $P$, il en va de même de leur
intersection $P'$. Celle-ci est stable sous $\pi' = \bigcap_{i} \pi_{i}$, et
ses fibres géométriques sont non vides puisqu'elles contiennent celle de
$P^{0}$. Donc \emph{$P'$ est un fibré principal homogène sous $\pi'$}.

\subsection*{La conséquence}

\textbf{Conséquence.} Si $\pi$ varie dans une catégorie de groupes artiniens
stable par produits finis et par sous-groupes, le foncteur
\[
  \pi \longmapsto \pi^{1}(S, \xi\,;\, \pi)
\]
est \emph{strictement pro-représentable}.

\subsection*{L'application : ce que le groupe élargi classe}

La page 168 dit à quoi cela sert, et c'est l'énoncé qui donne son sens au lot.

Une relation étale est dite \emph{à groupe structural $\pi$} si le schéma des
isomorphismes $\underline{\mathrm{Isom}}_{S}(E_{S}, T)$ est un revêtement
principal de groupe $\pi_{S}$. Si $S$ est connexe muni d'un point géométrique
$\xi$, alors \emph{se donner un $T$ muni d'une telle structure revient à se
donner des opérations continues de $\tilde{\pi}_{1}(S, \xi)$ sur $F$},
c'est-à-dire un homomorphisme continu $\tilde{\pi}_{1}(S,\xi) \to \pi$.

\textbf{Exemple.} Les schémas en groupes commutatifs sans torsion sur $S$, de
type fini, correspondent aux $\tilde{\pi}_{1}(S,\xi)$-modules qui sont des
$\mathbb{Z}$-modules de type fini.

C'est exactement ce que le groupe fondamental ordinaire ne peut pas faire : un
$\mathbb{Z}$-module de type fini n'est pas un ensemble fini, et le groupe
profini ne classe que des ensembles finis. Le groupe élargi est la réponse à ce
manque.\footnote{Cet objet --- le groupe fondamental classifiant les objets
localement constants à fibres non nécessairement finies --- porte aujourd'hui
plusieurs noms selon la classe de coefficients qu'on autorise : \emph{groupe
fondamental élargi} dans SGA 3, et \emph{groupe fondamental pro-étale} dans le
travail de Bhatt et Scholze (2015), qui en donne la construction topologique
correcte pour un schéma quelconque. Le carton en donne ici la construction par
pro-représentabilité, et ne discute pas la topologie du pro-groupe obtenu.}

\pagerange{170}{171}
\section*{Les cinq classes de recouvrements (pages 170 et 171)}

Ces deux pages sont écrites au net et se lisent presque entièrement --- ce qui,
dans ce carton, mérite d'être signalé.

Soit $S$ un préschéma. Pour tout groupe \emph{discret} $G$ on pose
\[
  H^{1}_{\mathcal{C}}(S, G) \;=\; \varinjlim_{T/S} H^{1}(T/S, G) ,
  \qquad
  H^{1}(T/S, G) \;=\; H^{1}\bigl(\pi_{0}(K_{T/S}),\, G\bigr) ,
\]
la limite étant prise sur une catégorie $\mathcal{C}$ de $S$-préschémas. Les
cinq classes qui importent sont :

\begin{itemize}
  \item[$C_{1}$] les $S$-préschémas fidèlement plats quasi-compacts ;
  \item[$C_{2}$] les mêmes, de type fini ;
  \item[$C_{3}$] les mêmes, quasi-finis ;
  \item[$C_{4}$] les mêmes, finis principaux ;
  \item[$C_{5}$] les mêmes, finis étales principaux.
\end{itemize}

On a des applications injectives canoniques $H^{1}_{C_{i}}(S,G)
\hookrightarrow H^{1}_{C_{j}}(S,G)$ pour $C_{i} \subset C_{j}$.

\subsection*{Pourquoi cela n'a d'intérêt que pour $G$ infini}

Si $G$ est \emph{fini}, tous ces homomorphismes sont bijectifs, et
$H^{1}(S,G)$ est simplement l'ensemble des classes de revêtements étales
principaux de groupe $G$, soit
\[
  \operatorname{Hom}\bigl(\pi_{1}(S,a),\, G\bigr) \big/
  \operatorname{int}(G) .
\]
La raison est le théorème de descente : $G_{T}$ est alors fini, donc affine sur
$T$, et la descente fidèlement plate quasi-compacte des schémas affines est
effective. Toutes les classes donnent la même réponse, et l'appareil est vide.

D'où la phrase qui justifie le lot : \emph{« Notre intérêt ici est le cas $G$
infini »}.

\subsection*{Les comparaisons}

\begin{itemize}
  \item En général, $H^{1}_{C_{3}} \simeq H^{1}_{C_{2}}$ : quasi-fini ou de
    type fini, cela ne change rien.
  \item Si $S$ est le spectre d'un anneau de Dedekind,
    $H^{1}_{C_{2}} \xrightarrow{\ \sim\ } H^{1}_{C_{1}}$, d'où
    $H^{1}_{C_{3}} \simeq H^{1}_{C_{2}} \simeq H^{1}_{C_{1}}$.
  \item Si $S$ est le spectre d'un anneau de valuation discrète complet à corps
    résiduel algébriquement clos, $H^{1}_{C_{i}} \simeq H^{1}(\hat{V}/V, -)$
    pour $i = 1, 2, 3$ : une seule extension suffit.
  \item Si $S$ est le spectre d'un anneau local \emph{complet} $V$ de corps
    résiduel $k$, alors $H^{1}_{C_{5}} \simeq H^{1}_{C_{4}} \simeq
    H^{1}_{C_{3}} \simeq H^{1}_{C_{2}}$, et de plus
    $H^{1}_{C_{i}}(V, -) \simeq H^{1}_{C_{i}}(k, -)$ pour $i = 2,3,4,5$.
\end{itemize}

La dernière ligne est la version « groupes infinis » du fait que le groupe
fondamental d'un anneau local complet est celui de son corps résiduel.

\pagerange{173}{173}
\section*{Un plan en sept sections (page 173)}

Le feuillet qui précède porte, de sa main, le seul mot « Plan ». Voici ce qu'il
annonce, pour deux théorèmes de changement de base :

\begin{itemize}
  \item[I] Premier théorème de changement de base pour $\pi_{1}$, ou théorème
    de comparaison.
  \item[II] Deuxième théorème de changement de base pour $\pi_{1}$.
\end{itemize}

\begin{enumerate}
  \item Morphismes $(-1)$-connexes et $0$-connexes.
  \item Morphismes \emph{locaux} $(-1)$-connexes et $0$-connexes.
  \item Caractérisation des morphismes locaux ou globaux $i$-connexes
    ($i = 0, 1$).
  \item Deuxième théorème du changement de base pour $\pi_{1}$ : énoncés et
    corollaires.
  \item Démonstration des théorèmes.
  \item Généralisation à l'aide de la résolution des singularités.
  \item Application : la dimension cohomologique de certains corps.
\end{enumerate}

Le vocabulaire des morphismes $i$-connexes --- un morphisme est $(-1)$-connexe
s'il est surjectif au sens des faisceaux, $0$-connexe s'il induit de plus une
bijection sur les $\pi_{0}$, $1$-connexe s'il induit un isomorphisme sur les
$\pi_{1}$ --- est celui de l'homotopie, transposé. Le plan n'est pas exécuté
dans ce carton.

\pagerange{175}{176}
\section*{$\pi_{2}$ contre la tour des revêtements (pages 175 et 176)}

Le feuillet qui ouvre la suite porte, de sa main, « $\Pi_{2}$ ». Ce sont deux
pages de calcul, sans phrases.

\subsection*{Le calcul}

On dispose de la tour d'un revêtement $X^{i} \to X$, de son revêtement
universel $\tilde{X}$, et des revêtements universels itérés
$\tilde{\tilde{X}}^{i}$ :

\begin{tikzcd}[column sep=small, row sep=small, nodes={font=\scriptsize}]
\widetilde{\widetilde{X}}{}^i \arrow[d] & \\
\widetilde{X}^i \arrow[d] & \widetilde{X} \arrow[d, no head] \\
X^i \arrow[r] & X
\end{tikzcd}

et l'on pose $\tilde{\pi}/H_{i} = \pi_{i}$. Le calcul est
\[
  \pi^{2}(X, G) \;\simeq\; H^{2}(\check{X}, G) \;\simeq\;
  \varinjlim_{i} H^{2}(X_{i}, G) \;\simeq\;
  \varinjlim_{i} H^{2}(\tilde{X}_{i}, G) ,
\]
la suite spectrale de Hochschild--Serre donnant
\[
  E_{2}^{p,q} \;=\; H^{p}\bigl(H_{i},\, H^{q}(\tilde{X}, G)\bigr)
  \;\Longrightarrow\; H^{*}(\tilde{X}_{i}, G) ,
\]
puis, en passant à la limite sur $i$, la suite exacte
\[
  \cdots \to \varinjlim_{i} H^{2}(H_{i}, G) \to \pi^{2}(X, G) \to
  \varinjlim_{i} \pi^{2}(\tilde{X}_{i}, G)^{H_{i}} \to
  \varinjlim_{i} H^{3}(H_{i}, G) .
\]
Grothendieck entoure les deux extrémités, qui sont ce qu'il voulait :
\[
  \operatorname{Ker}\bigl(\pi^{2}(X,G) \to \pi^{2}(X^{k},G)\bigr) \;\simeq\;
  \varinjlim_{i} H^{2}(H_{i}, G) ,
\]
\[
  \operatorname{Coker} \;\simeq\;
  \operatorname{Ker}\Bigl(\varinjlim_{i} H^{3}(H_{i}, G) \longrightarrow
  H^{3}(\tilde{X}, G)\Bigr) .
\]
Autrement dit : \emph{ce que le passage aux revêtements fait perdre et gagne en
$\pi_{2}$ est de la cohomologie des groupes de la tour}.\footnote{L'exposant du
premier terme est tracé comme un « I » romain sur la page, là où quatre lignes
plus bas le même objet est noté $\pi^{2}$ ; la transcription le signale et le
conserve. C'est $\pi^{2}$ qui est écrit ici.}

\subsection*{Le cas d'une fibration}

La page 176 fait la même chose pour une fibration $\bar{F} \to X \to Y$, avec
la suite exacte d'homotopie $\pi_{1}(\bar{F}) \to \pi_{1}(X) \to \pi_{1}(Y)
\to e$, et aboutit à l'annulation
\[
  H^{1}\bigl(Y'/Y,\ \operatorname{Hom}(\pi_{1}(\bar{F}),\, \mathfrak{z})\bigr)
  \;=\; 0 ,
\]
qui est la condition sous laquelle le passage de la fibre à l'espace total est
sans obstruction.\footnote{Le groupe noté $\mathfrak{z}$ ne porte pas de nom
sur la page ; c'est le centralisateur $G^{\pi_{1}(\bar{F})}$ qui apparaît deux
lignes plus haut. La page est faite de formules sans phrases liantes, et cette
lecture n'invente pas les phrases : elle donne le calcul et son point
d'arrivée.}

\pagerange{178}{180}
\section*{Descente : pro-représentabilité, et une question ouverte (pages 178 à 180)}

Le feuillet qui ouvre la suite porte, de sa main : « $\pi_{0}$, $\pi_{1}$ et
formalisme simplicial (Revêtements infinis) ».

\subsection*{La proposition}

Soit $S$ un préschéma \emph{connexe}, $S' \to S$ un morphisme
\emph{submersif} --- par exemple fidèlement plat quasi-compact, ou propre
surjectif --- et $S''$, $S'''$ le carré et le cube fibrés. On suppose que $S'$
est \emph{somme} de préschémas connexes, c'est-à-dire que $\pi_{0}(S')$ est
discret. Soit $\xi$ un point géométrique de $S$.

Alors, pour un groupe $G$ variable, le foncteur \emph{covariant}
\[
  G \longmapsto \pi^{1}(S'/S, \xi\,;\, G)
\]
--- l'ensemble des classes, à isomorphisme près, de \emph{données de descente}
sur $G_{S'}$ relatives à $S' \to S$, pointées au-dessus de $\xi$ --- est
strictement pro-représentable :
\[
  \pi^{1}(S'/S, \xi\,;\, G) \;=\; \varinjlim_{i}
  \operatorname{Hom}(\pi_{i},\, G) .
\]

Deux raffinements. Si $\pi_{0}(S'')$ est quasi-compact, les $\pi_{i}$ sont de
type fini. Et si $S''$ n'a qu'un nombre fini $n''$ de composantes connexes, le
foncteur est \emph{représentable} :
\[
  \pi^{1}(S'/S, \xi\,;\, G) \;\simeq\;
  \operatorname{Hom}\bigl(\pi_{1}(S'/S, \xi),\, G\bigr) ,
\]
avec $\pi_{1}(S'/S,\xi)$ un groupe \emph{discret} de type fini, engendré par
$n''$ générateurs.\footnote{La page écrit d'abord « contravariant », puis le
biffe. C'est bien un foncteur covariant : une donnée de descente sur $G_{S'}$
se pousse le long d'un homomorphisme $G \to G'$. La correction est de sa main
et elle est suivie ici.}

C'est le \emph{groupe fondamental de la descente} : un groupe discret de type
fini, engendré par autant d'éléments que $S''$ a de composantes connexes. La
finitude du nombre de générateurs vient donc de la combinatoire du
recouvrement, non de la géométrie --- exactement comme au contre-exemple de la
page 161.

\subsection*{La question ouverte}

Considérons de même le foncteur
$G \mapsto \pi^{1}_{\mathrm{eff}}(S'/S, \xi\,;\, G)$ des données de descente
\emph{effectives} pointées. Il est également pro-représentable, par un système
projectif de \emph{quotients} des $\pi_{i}$.

Et vient l'aveu :

\begin{quote}
J'ignore, si $S$, $S'$, $S''$ sont noethériens, si ce dernier foncteur est
représentable ; \emph{a fortiori} j'ignore si dans ce cas
$\pi^{1} = \pi^{1}_{\mathrm{eff}}$, c'est-à-dire si \emph{toute} donnée de
descente sur $G_{S'}$ est \emph{effective}.
\end{quote}

Il ajoute que c'est vrai sous une hypothèse supplémentaire, et note en marge :
« Pb intér.\ pour l'affirmative [\dots\ et $S' \to S$ fidèlement plat
quasi-compact] ».

Il vaut la peine de dire pourquoi la question est ouverte, parce que la raison
se lit en rapprochant cette page de la page 170.

Lorsque $G$ est \emph{fini}, le schéma en groupes constant $G_{S'}$ est fini
sur $S'$, donc \emph{affine} sur $S'$ ; et le théorème de descente fidèlement
plate de Grothendieck dit précisément que la descente est effective pour les
morphismes affines. La question est alors réglée, et c'est ce que la page 170
constatait en disant que les cinq classes donnent la même réponse pour $G$
fini.

Lorsque $G$ est \emph{infini}, $G_{S'}$ est la somme disjointe d'une famille
infinie de copies de $S'$ : elle n'est pas quasi-compacte sur $S'$, donc pas
affine, et le théorème de descente ne l'atteint plus. \emph{La question est
donc ouverte exactement dans le cas que la page 170 déclarait être le seul
intéressant.}

\subsection*{La traduction simpliciale}

La page 180 réduit tout cela à de la combinatoire. On dispose du diagramme
semi-simplicial des $X_{i}$, de leurs fibres $X_{i}^{\xi}$, et des ensembles
\[
  K_{i}^{\xi} \;=\; \pi_{0}(X_{i}^{\xi}) ,
\]
avec, en dessous, la version pointée $k_{i}$ et les inclusions $i_{j}$. Alors :

\begin{quote}
Les données de descente sur $E_{X_{0}}$ correspondent aux données de descente
sur $E_{K_{0}}$, c'est-à-dire aux applications ensemblistes
$f : K_{1} \to G$ vérifiant
\[
  f \circ p_{1} \;=\; (f \circ p_{0})\,(f \circ p_{2}) ,
\]
deux telles applications étant équivalentes si
$f'(q) = h(p_{1}q)\, f(q)\, h(p_{0}q)^{-1}$ ; et les données \emph{pointées}
sont celles qui vérifient de plus $f(i_{1}(q)) = e_{G}$ pour tout
$q \in k_{1}$, l'équivalence étant restreinte aux $h$ tels que
$h(i_{0}(q)) = 1$.
\end{quote}

C'est la relation de cocycle non abélienne en degré 1, et c'est la même que
celle que la page 144 encadrait : $A_{\alpha''} a''_{\beta} A_{\alpha'} =
a'''_{\beta} A_{\alpha'''} a'_{\beta}$. Elle est ici écrite proprement, sur
un ensemble semi-simplicial, ce qui rend visible sa nature : \emph{un
$1$-cocycle non abélien sur le nerf du recouvrement}. La
pro-représentabilité de la proposition ci-dessus n'est alors rien d'autre que
la présentation par générateurs de ce nerf --- d'où les $n''$ générateurs. Le
lot 10 poursuit ce calcul, sur des foncteurs semi-cosimpliciaux abstraits.

\end{document}
